\documentclass[a4paper,12pt]{report}
\usepackage[utf8]{inputenc}
\usepackage[T1]{fontenc}
\usepackage[french]{babel}
\usepackage{geometry}
\geometry{margin=2.5cm}
\usepackage{graphicx}
\usepackage{enumitem}
\usepackage{array}
\usepackage{booktabs}
\usepackage{tabularx}
\usepackage{xcolor}
\usepackage{titlesec}
\usepackage{fancyhdr}
\usepackage{hyperref}
\usepackage{tikz}
\usepackage{tcolorbox}
\usepackage{pgfgantt}
\usepackage{amssymb}

% Couleurs BPM
\definecolor{bpmred}{HTML}{C8102E}
\definecolor{bpmgreen}{HTML}{00703C}
\definecolor{bpmblue}{HTML}{1A3A5C}
\definecolor{bpmgray}{HTML}{F5F5F5}
\definecolor{done}{HTML}{27AE60}
\definecolor{inprogress}{HTML}{F39C12}
\definecolor{todo}{HTML}{E74C3C}

% Titres
\titleformat{\chapter}[display]{\normalfont\Huge\bfseries\color{bpmred}}{\chaptertitlename\ \thechapter}{20pt}{\Huge}
\titleformat{\section}{\normalfont\Large\bfseries\color{bpmblue}}{\thesection}{1em}{}
\titleformat{\subsection}{\normalfont\large\bfseries\color{bpmblue!80}}{\thesubsection}{1em}{}

% En-tête/pied de page
\pagestyle{fancy}
\fancyhf{}
\fancyhead[L]{\small\textcolor{bpmblue}{État d'Avancement — GAB Intelligence Dashboard}}
\fancyhead[R]{\small\textcolor{bpmblue}{Banque Populaire du Maroc}}
\fancyfoot[C]{\thepage}
\renewcommand{\headrulewidth}{0.4pt}

\hypersetup{
    colorlinks=true,
    linkcolor=bpmblue,
    urlcolor=bpmred,
}

% Commandes
\newcommand{\done}{\textcolor{done}{\checkmark\ Terminé}}
\newcommand{\inprog}{\textcolor{inprogress}{$\circlearrowright$\ En cours}}
\newcommand{\notstarted}{\textcolor{todo}{$\circ$\ Non démarré}}

\begin{document}

% ===== PAGE DE GARDE =====
\begin{titlepage}
\centering
\vspace*{1cm}

{\Large\textsc{Banque Populaire du Maroc}}\par
\vspace{0.5cm}
{\large Projet de Fin d'Études — 2024/2025}\par

\vspace{2cm}
\rule{\textwidth}{2pt}\par
\vspace{0.8cm}
{\Huge\bfseries\textcolor{bpmred}{État d'Avancement}}\par
\vspace{0.5cm}
{\LARGE\textcolor{bpmblue}{Système de Détection et de Prédiction\\des Pannes des GAB}}\par
\vspace{0.3cm}
{\Large GAB Intelligence Dashboard}\par
\vspace{0.8cm}
\rule{\textwidth}{2pt}\par

\vspace{1.5cm}

\begin{tabular}{ll}
\textbf{Organisme d'accueil :} & Banque Populaire du Maroc \\
\textbf{Encadrant entreprise :} & \textit{[Nom de l'encadrant]} \\
\textbf{Encadrant académique :} & \textit{[Nom de l'encadrant]} \\
\textbf{Réalisé par :} & \textit{[Votre nom complet]} \\
\textbf{Date du rapport :} & \today \\
\end{tabular}

\vspace{2cm}

\begin{tcolorbox}[colback=done!10, colframe=done, title=\textbf{Avancement global}]
\centering
{\Huge\textcolor{done}{\textbf{95\%}}}\\[0.5em]
{\large Le système est fonctionnel et opérationnel.\\Reste : tests finaux et rédaction du rapport de soutenance.}
\end{tcolorbox}

\vfill
{\small Document confidentiel — Diffusion restreinte}
\end{titlepage}

% ===== TABLE DES MATIÈRES =====
\tableofcontents
\newpage

% ===================================================================
\chapter{Résumé Exécutif}
% ===================================================================

Le projet \textbf{GAB Intelligence Dashboard} est un système de détection et de prédiction de pannes des Guichets Automatiques Bancaires de la Banque Populaire du Maroc. Il combine des techniques de \textbf{Machine Learning} avec un \textbf{tableau de bord interactif} pour permettre aux équipes de maintenance de passer d'une approche corrective à une approche prédictive.

\begin{tcolorbox}[colback=bpmgray, colframe=bpmblue, title=\textbf{Chiffres clés}]
\begin{itemize}[leftmargin=1cm]
    \item \textbf{200 GAB} surveillés dans \textbf{13 villes} marocaines
    \item \textbf{146\,000 observations} analysées sur 2 ans (2022–2023)
    \item \textbf{101 features} d'ingénierie extraites en 7 familles
    \item \textbf{5 modèles ML} entraînés, évalués et comparés
    \item \textbf{Prédiction 48h à l'avance} avec un F1-score $\geq$ 0.70
    \item \textbf{6 modules} fonctionnels entièrement développés
    \item \textbf{Économies potentielles $>$ 50\%} sur les coûts de maintenance
\end{itemize}
\end{tcolorbox}

% ===================================================================
\chapter{Avancement par Phase}
% ===================================================================

\section{Vue synthétique}

\begin{table}[h!]
\centering
\renewcommand{\arraystretch}{1.5}
\begin{tabularx}{\textwidth}{|c|X|c|c|}
\hline
\rowcolor{bpmblue!10}
\textbf{\#} & \textbf{Phase} & \textbf{Avancement} & \textbf{Statut} \\
\hline
1 & Étude de l'existant et analyse des besoins & 100\% & \done \\
\hline
2 & Collecte et préparation des données & 100\% & \done \\
\hline
3 & Ingénierie des features et entraînement ML & 100\% & \done \\
\hline
4 & Développement du backend (API REST) & 100\% & \done \\
\hline
5 & Développement du frontend (Dashboard) & 100\% & \done \\
\hline
6 & Intégration, tests et optimisation & 90\% & \inprog \\
\hline
7 & Rédaction du rapport et soutenance & 60\% & \inprog \\
\hline
\end{tabularx}
\caption{Avancement global par phase}
\end{table}

% ===================================================================
\section{Phase 1 — Étude de l'existant et analyse des besoins}

\textbf{Statut : \done}

\begin{itemize}
    \item Étude du parc GAB de la Banque Populaire (200 automates, 4 marques, 13 villes).
    \item Analyse des processus de maintenance existants (modèle correctif).
    \item Identification des indicateurs clés~: taux de panne, MTBF, coûts d'intervention.
    \item Rédaction du cahier des charges fonctionnel et technique.
    \item Choix de l'architecture technique (FastAPI + React + scikit-learn).
\end{itemize}

% ===================================================================
\section{Phase 2 — Collecte et préparation des données}

\textbf{Statut : \done}

\begin{itemize}
    \item Constitution du jeu de données~: \textbf{146\,000 observations} sur la période janvier 2022 – décembre 2023.
    \item Variables collectées~: identifiant GAB, ville, type, environnement, âge, nombre d'erreurs, température, jours depuis maintenance, nombre de transactions, etc.
    \item Nettoyage et traitement des valeurs manquantes.
    \item Analyse exploratoire des données (EDA)~:
    \begin{itemize}
        \item Taux de panne global~: \textbf{9,8\%} (classe déséquilibrée).
        \item Villes les plus affectées~: Safi, Agadir, Casablanca ($\approx$ 10–11\%).
        \item Variations saisonnières significatives identifiées.
    \end{itemize}
\end{itemize}

% ===================================================================
\section{Phase 3 — Ingénierie des features et entraînement ML}

\textbf{Statut : \done}

\subsection{Ingénierie des features}

Construction de \textbf{101 features} réparties en 7 familles~:

\begin{table}[h!]
\centering
\renewcommand{\arraystretch}{1.2}
\begin{tabular}{|l|l|c|}
\hline
\rowcolor{bpmblue!10}
\textbf{Famille} & \textbf{Exemples} & \textbf{Nombre} \\
\hline
Original & âge, erreurs, température, transactions & $\sim$15 \\
\hline
Lag (J-1/3/7) & erreurs\_j1, temperature\_j3, transactions\_j7 & $\sim$20 \\
\hline
Rolling (7-14j) & erreurs\_moy\_7j, temp\_std\_14j & $\sim$20 \\
\hline
Tendance & pente\_erreurs, delta\_temperature & $\sim$10 \\
\hline
Temporel & mois, jour\_semaine, saison & $\sim$10 \\
\hline
Catégoriel & ville\_encoded, type\_gab\_encoded & $\sim$15 \\
\hline
Interaction & age $\times$ erreurs, maintenance $\times$ temp & $\sim$11 \\
\hline
\end{tabular}
\caption{Familles de features construites}
\end{table}

\subsection{Entraînement et évaluation}

\begin{itemize}
    \item \textbf{Découpage temporel}~: entraînement sur 2022, test sur 2023.
    \item \textbf{5 modèles} entraînés et comparés~:
\end{itemize}

\begin{table}[h!]
\centering
\renewcommand{\arraystretch}{1.3}
\begin{tabular}{|l|c|c|c|c|c|}
\hline
\rowcolor{bpmblue!10}
\textbf{Modèle} & \textbf{F1} & \textbf{Précision} & \textbf{Rappel} & \textbf{AUC-ROC} & \textbf{AUC-PR} \\
\hline
Régression Logistique & $\sim$0.55 & $\sim$0.60 & $\sim$0.52 & $\sim$0.80 & $\sim$0.40 \\
\hline
\textbf{Random Forest} & $\sim$\textbf{0.72} & $\sim$\textbf{0.68} & $\sim$\textbf{0.76} & $\sim$\textbf{0.92} & $\sim$\textbf{0.65} \\
\hline
Gradient Boosting & $\sim$0.70 & $\sim$0.66 & $\sim$0.74 & $\sim$0.91 & $\sim$0.62 \\
\hline
SVM & $\sim$0.58 & $\sim$0.62 & $\sim$0.55 & $\sim$0.82 & $\sim$0.42 \\
\hline
Réseau de Neurones & $\sim$0.65 & $\sim$0.63 & $\sim$0.67 & $\sim$0.87 & $\sim$0.55 \\
\hline
\textit{Dummy (baseline)} & $\sim$0.17 & $\sim$0.10 & $\sim$0.50 & $\sim$0.50 & $\sim$0.10 \\
\hline
\end{tabular}
\caption{Résultats comparatifs des modèles (valeurs approximatives)}
\end{table}

\begin{itemize}
    \item \textbf{Meilleur modèle}~: Random Forest avec calibration isotonique.
    \item \textbf{Calibration}~: régression isotonique appliquée pour des probabilités fiables.
    \item Sérialisation du modèle et des encodeurs en fichiers \texttt{.pkl}.
\end{itemize}

\textit{Note~: Les valeurs exactes sont disponibles dans le fichier \texttt{resultats\_modeles.json} du projet.}

% ===================================================================
\section{Phase 4 — Développement du backend (API REST)}

\textbf{Statut : \done}

\subsection{Endpoints développés}

\textbf{12 endpoints} API REST implémentés avec FastAPI~:

\begin{itemize}
    \item \texttt{GET /api/health} — État du système et mode du modèle (ML ou heuristique).
    \item \texttt{GET /api/metadata} — Métadonnées (villes, types GAB, années disponibles).
    \item \texttt{GET /api/overview} — KPIs globaux avec filtres multi-critères.
    \item \texttt{GET /api/geography} — Analyse géographique et heatmap villes $\times$ mois.
    \item \texttt{GET /api/models} — Métriques comparatives de tous les modèles ML.
    \item \texttt{GET /api/features} — Importance des variables avec classement par famille.
    \item \texttt{GET /api/threshold} — Simulation économique de seuils de décision.
    \item \texttt{POST /api/scoring} — Scoring en temps réel d'un GAB individuel.
    \item \texttt{GET /api/monitoring} — Surveillance temps réel de l'ensemble du parc.
    \item \texttt{GET /api/predictions-stats} — Statistiques du journal d'audit.
    \item \texttt{POST /api/reload-model} — Rechargement à chaud du modèle (admin).
    \item \texttt{GET /api/export/*} — Export des données en CSV, Excel et PDF.
\end{itemize}

\subsection{Fonctionnalités techniques}

\begin{itemize}
    \item \textbf{Cache intelligent} — TTL configurable (300s par défaut, 25s pour le monitoring).
    \item \textbf{Mode dégradé} — Algorithme heuristique avec 9 composantes pondérées.
    \item \textbf{Journal d'audit} — Toutes les prédictions enregistrées en JSONL.
    \item \textbf{CORS} — Configuration pour le développement local.
\end{itemize}

% ===================================================================
\section{Phase 5 — Développement du frontend (Dashboard)}

\textbf{Statut : \done}

\subsection{Pages développées}

\begin{table}[h!]
\centering
\renewcommand{\arraystretch}{1.4}
\begin{tabularx}{\textwidth}{|l|l|X|}
\hline
\rowcolor{bpmblue!10}
\textbf{Page} & \textbf{Route} & \textbf{Contenu} \\
\hline
Vue d'ensemble & \texttt{/} & 4 KPIs + 5 graphiques (tendances, types, environnements, saisons, âge) \\
\hline
Géographie & \texttt{/geo} & Classement des villes + heatmap interactive \\
\hline
Modèles & \texttt{/modeles} & Radar comparatif + matrice de confusion \\
\hline
Features & \texttt{/features} & Bar chart importance + répartition par famille \\
\hline
Seuil & \texttt{/seuil} & Simulation coûts/bénéfices + seuil optimal \\
\hline
Scoring & \texttt{/scoring} & 15 paramètres + jauge de risque + contributions \\
\hline
\end{tabularx}
\caption{Pages du tableau de bord}
\end{table}

\subsection{Composants réutilisables}

\begin{itemize}
    \item \texttt{useApi} / \texttt{usePost} — Hooks de communication avec le backend.
    \item \texttt{Skeleton} / \texttt{PageSkeleton} — Chargement avec animation shimmer.
    \item \texttt{ErrorState} — Gestion unifiée des erreurs.
    \item \texttt{Chart} — Wrapper Plotly avec palette et layout standardisés.
    \item \texttt{Icons} — Bibliothèque d'icônes SVG et composants visuels.
    \item \texttt{FilterContext} — Filtres globaux partagés entre toutes les pages.
\end{itemize}

\subsection{Fonctionnalités UX}

\begin{itemize}
    \item Thème clair/sombre avec persistance localStorage.
    \item Graphiques interactifs (zoom, hover, sélection).
    \item Alertes temps réel avec icône de notification (cloche).
    \item Export multi-format (CSV, Excel, PDF).
    \item Barre latérale de navigation avec logo BPM.
\end{itemize}

% ===================================================================
\section{Phase 6 — Intégration, tests et optimisation}

\textbf{Statut : \inprog\ (90\%)}

\begin{table}[h!]
\centering
\renewcommand{\arraystretch}{1.3}
\begin{tabularx}{\textwidth}{|X|c|}
\hline
\rowcolor{bpmblue!10}
\textbf{Tâche} & \textbf{Statut} \\
\hline
Intégration frontend-backend complète & \done \\
\hline
Tests de bout en bout de tous les modules & \done \\
\hline
Script de démarrage automatisé (\texttt{start.bat}) & \done \\
\hline
Optimisation des performances (cache, lazy loading) & \done \\
\hline
Export PDF avec mise en forme professionnelle & \done \\
\hline
Tests de charge et de stress & \inprog \\
\hline
Tests sur différents navigateurs & \inprog \\
\hline
\end{tabularx}
\caption{Détail de la phase d'intégration}
\end{table}

% ===================================================================
\section{Phase 7 — Rédaction du rapport et soutenance}

\textbf{Statut : \inprog\ (60\%)}

\begin{table}[h!]
\centering
\renewcommand{\arraystretch}{1.3}
\begin{tabularx}{\textwidth}{|X|c|}
\hline
\rowcolor{bpmblue!10}
\textbf{Tâche} & \textbf{Statut} \\
\hline
Rédaction du cahier des charges & \done \\
\hline
Rédaction de l'état d'avancement & \done \\
\hline
Rédaction du rapport de PFE & \inprog \\
\hline
Préparation de la présentation de soutenance & \notstarted \\
\hline
\end{tabularx}
\caption{Détail de la phase de rédaction}
\end{table}

% ===================================================================
\chapter{Démonstration du Système}
% ===================================================================

\section{Captures d'écran}

\begin{figure}[h!]
    \centering
    \includegraphics[width=0.95\textwidth]{../outputs/figures/fig1_vue_ensemble.png}
    \caption{Vue d'ensemble du jeu de données — statistiques descriptives et répartition des variables}
    \label{fig:vue_ensemble}
\end{figure}

\begin{figure}[h!]
    \centering
    \includegraphics[width=0.95\textwidth]{../outputs/figures/fig2_distributions.png}
    \caption{Distributions des variables numériques clés}
    \label{fig:distributions}
\end{figure}

\begin{figure}[h!]
    \centering
    \includegraphics[width=0.95\textwidth]{../outputs/figures/fig3_correlations.png}
    \caption{Matrice de corrélation entre les variables principales}
    \label{fig:correlations}
\end{figure}

\begin{figure}[h!]
    \centering
    \includegraphics[width=0.95\textwidth]{../outputs/figures/fig4_boxplots_top6.png}
    \caption{Boxplots des 6 variables les plus discriminantes}
    \label{fig:boxplots}
\end{figure}

\begin{figure}[h!]
    \centering
    \includegraphics[width=0.95\textwidth]{../outputs/figures/fig5_analyse_temporelle.png}
    \caption{Analyse temporelle — évolution des pannes et tendances saisonnières}
    \label{fig:temporelle}
\end{figure}

\begin{figure}[h!]
    \centering
    \includegraphics[width=0.95\textwidth]{../outputs/models/fig7_comparaison_modeles.png}
    \caption{Comparaison des performances des 5 modèles ML}
    \label{fig:comparaison}
\end{figure}

\begin{figure}[h!]
    \centering
    \includegraphics[width=0.95\textwidth]{../outputs/models/fig8_roc_pr_cm.png}
    \caption{Courbes ROC, Precision-Recall et matrice de confusion du meilleur modèle}
    \label{fig:roc_pr_cm}
\end{figure}

\begin{figure}[h!]
    \centering
    \includegraphics[width=0.95\textwidth]{../outputs/models/fig9_analyse_metier.png}
    \caption{Analyse métier — impact économique et optimisation du seuil de décision}
    \label{fig:analyse_metier}
\end{figure}

\section{Accès au système}

\begin{tcolorbox}[colback=bpmgray, colframe=bpmblue, title=\textbf{Lancement du système}]
\begin{enumerate}
    \item Exécuter le script \texttt{start.bat} à la racine du projet.
    \item Le backend FastAPI démarre sur \texttt{http://localhost:8000}.
    \item Le frontend React démarre sur \texttt{http://localhost:5173}.
    \item Le navigateur s'ouvre automatiquement sur le tableau de bord.
\end{enumerate}
\end{tcolorbox}

% ===================================================================
\chapter{Résultats et Impact}
% ===================================================================

\section{Résultats techniques}

\begin{itemize}
    \item Le meilleur modèle (Random Forest) atteint un \textbf{F1-score $\geq$ 0.70} et un \textbf{rappel $\geq$ 0.75}, détectant plus de 3 pannes sur 4 avant qu'elles ne surviennent.
    \item L'analyse de seuil économique montre des \textbf{économies potentielles supérieures à 50\%} par rapport à une stratégie purement corrective.
    \item Le système de scoring temps réel fournit des prédictions en \textbf{moins de 200~ms}.
\end{itemize}

\section{Impact métier}

\begin{table}[h!]
\centering
\renewcommand{\arraystretch}{1.4}
\begin{tabularx}{\textwidth}{|l|X|}
\hline
\rowcolor{bpmblue!10}
\textbf{Indicateur} & \textbf{Amélioration attendue} \\
\hline
Coût de maintenance & Réduction $>$ 50\% (passage correctif $\rightarrow$ préventif) \\
\hline
Disponibilité des GAB & Augmentation significative (prédiction 48h à l'avance) \\
\hline
Satisfaction client & Amélioration (moins d'automates hors service) \\
\hline
Temps de décision & Réduction (tableau de bord temps réel vs rapports manuels) \\
\hline
Traçabilité & 100\% des prédictions journalisées \\
\hline
\end{tabularx}
\caption{Impact métier attendu}
\end{table}

% ===================================================================
\chapter{Difficultés Rencontrées et Solutions}
% ===================================================================

\begin{table}[h!]
\centering
\renewcommand{\arraystretch}{1.4}
\begin{tabularx}{\textwidth}{|X|X|}
\hline
\rowcolor{bpmblue!10}
\textbf{Difficulté} & \textbf{Solution apportée} \\
\hline
Déséquilibre de classes (9,8\% positifs) & Métriques adaptées (F1, AUC-PR au lieu de l'accuracy) \\
\hline
Risque de fuite temporelle & Découpage temporel strict (train 2022 / test 2023) \\
\hline
Taille du bundle frontend (Plotly.js $>$ 5 MB) & Utilisation de \texttt{plotly.js-dist-min}, warning accepté \\
\hline
Indisponibilité potentielle du modèle ML & Implémentation d'un mode dégradé heuristique \\
\hline
Calibration des probabilités & Application de la régression isotonique \\
\hline
\end{tabularx}
\caption{Difficultés et solutions}
\end{table}

% ===================================================================
\chapter{Prochaines Étapes}
% ===================================================================

\begin{enumerate}
    \item \textbf{Finaliser les tests} — Compléter les tests de charge et de compatibilité navigateurs.
    \item \textbf{Rédiger le rapport de PFE} — Documenter l'ensemble de la démarche, les résultats et les perspectives.
    \item \textbf{Préparer la soutenance} — Créer la présentation et préparer la démonstration en direct.
\end{enumerate}

\subsection*{Perspectives d'évolution}

\begin{itemize}
    \item Déploiement sur un serveur de production (Docker, cloud).
    \item Intégration avec le système d'information de la BPM (données temps réel).
    \item Réentraînement périodique automatisé des modèles.
    \item Extension à d'autres types d'équipements bancaires.
    \item Ajout de notifications par email/SMS pour les alertes critiques.
\end{itemize}

\end{document}
